% Requirements and tradeoffs, rather than an overall ranking of zero-copy mechanisms.
% Keep deployment requirements distinct from access, retention and buffer-lifetime obligations.
\section{What Zero-Copy Interfaces Require}
\label{sec:borrowing}

Zero-copy interfaces grant different access and reuse, with requirements for protection,
buffer lifetime and layout. Table~\ref{tab:space} separates these capabilities from their
costs; protection alone does not explain every limitation.

\begin{table*}[t]
\centering
\footnotesize
\setlength{\tabcolsep}{4pt}
\begin{tabular}{@{}
  >{\raggedright\arraybackslash}p{0.18\textwidth}
  >{\raggedright\arraybackslash}p{0.25\textwidth}
  >{\raggedright\arraybackslash}p{0.25\textwidth}
  >{\raggedright\arraybackslash}p{0.27\textwidth}@{}}
\toprule
Mechanism & Access and reuse & Setup requirements & Remaining costs or limits \\
\midrule
\multicolumn{4}{@{}l}{\emph{Existing Linux interfaces}} \\
\msgzc, default \sendzc~\cite{msgzerocopy,sendzcdoc}
  & Application buffers; repeated sends
  & Socket or \iouring send API
  & Pin per send; preserve bytes until reuse notification \\
registered \sendzc~\cite{sendzcdoc}
  & Registered buffers; repeated sends
  & Register and pin a buffer pool
  & Pinning is amortized; reuse notifications remain \\
hardware \recvzc~\cite{zcrx,zcrxreg}
  & Writable payload area; retained receive slots
  & Header split, flow steering, queues; administrative registration
  & Return slots; holding reduces the receive pool \\
devmem TCP~\cite{devmem}
  & Received ranges in device buffers
  & Device-buffer binding, header split and flow steering
  & Device-oriented access; release received ranges \\
TCP page remapping~\cite{tcpmmap}
  & Read-only mappings of received pages
  & Eligible full, aligned payload pages
  & Copy or skip ineligible bytes; manage mappings \\
\splice, \code{tee}~\cite{spliceapi}
  & Unmapped pipe data; \code{tee} duplicates references
  & Linux pipe APIs
  & Bounded FIFO pipes, not independently named ranges \\
page cache + \sendfile~\cite{sendfileapi}
  & File offsets; repeated range sends
  & Linux file API
  & Requires file-backed data; does not itself capture TCP input \\
sockmap~\cite{sockmap}
  & Socket streams redirected by BPF
  & Load a BPF program and attach sockets
  & Redirect interface; no retained-range handle \\
\addlinespace
\multicolumn{4}{@{}l}{\emph{Research designs}} \\
IO-Lite~\cite{iolite}
  & Read-only shared buffers; retained slices
  & Modified OS buffering and drivers
  & Choose access rights before receive placement \\
PASTE~\cite{paste}
  & Mapped buffers; retention and reuse; kernel TCP option
  & netmap extension and driver support
  & Isolated applications need DMA partitioning; buffer slack \\
MAIO~\cite{maio}
  & Segregated mapped I/O pages; kernel TCP
  & Kernel memory and driver changes; hardware header split
  & Per-process queues or privileged sharing \\
\addlinespace
\multicolumn{4}{@{}l}{\emph{\sys kernel extensions}} \\
\textbf{Named mode}
  & Read-only by handle; copied reads; concurrent range sends
  & Kernel patches and application API changes
  & Fallback copies, unused buffer space and object metadata \\
\textbf{Mapped mode}
  & Read-write I/O area; raw buffers or object handles
  & Kernel/driver/application changes; administrative privilege
  & Device-wide trust; coordinated writes; bounded receive pool \\
\bottomrule
\end{tabular}
\caption{Representative interfaces that retain kernel TCP. Requirements describe the
copy-avoiding path; fallbacks may copy. Buffers, pipes and files all support reuse. A reuse
notification permits modification after all sends using the bytes release them, rather
than merely reporting operation completion. Named mode maps nothing; mapped mode keeps
kernel control state and parsing copies outside its read-write I/O area.}
\label{tab:space}
\end{table*}

\subsection{Protection, access and retention}
\label{sec:borrowing:services}
\label{sec:borrowing:gap}

Hardware \recvzc keeps parsed headers private and directs flows into separate receive
pools~\cite{zcrx}. MAIO similarly uses header split, with per-process queues or privileged
sharing~\cite{maio}. The latter accepts access to other flows rather than isolating them.

IO-Lite chooses access rights in software before receive placement; PASTE retains mapped
data and needs a dedicated port or DMA partitioning for isolated applications~\cite{iolite,paste}.
Page remapping exposes eligible pages after receive, with alignment restrictions~\cite{tcpmmap}.
Retention and isolation do not inherently require specialized NIC features.

Pipes and files also support reuse: \code{tee} duplicates pipe references, and \sendfile
accepts file ranges~\cite{spliceapi,sendfileapi}. They expose bounded FIFO queues or file-backed
data, rather than independently named ranges captured from TCP.

\subsection{Costs depend on the interface}
\label{sec:borrowing:costs}

\paragraph{Resident memory.} The kernel must keep shared pages resident while it uses them.
Default \sendzc and \msgzc pin application pages for each send; registering buffers moves
that work to registration~\cite{msgzerocopy,sendzcdoc}. Kernel-owned buffers avoid pinning
application pages, but still consume resident memory and need limits on retention.

\paragraph{Safe reuse.} An application must not modify bytes that a send still uses. A reuse
notification tells it when modification is safe; this need not stall every send, and
notifications can be batched~\cite{msgzerocopy,sendzcdoc}. On receive, returning buffers makes
room for later packets. Holding buffers is allowed, but reduces the free pool.

\paragraph{Receive placement.} Hardware \recvzc and device-memory TCP require separate
headers and steered queues~\cite{zcrx,devmem}. Queue leasing extends this arrangement to
containers but still needs physical queues and steering rules~\cite{netkit}. These costs
differ from software buffer selection or page remapping.

\paragraph{Stack handling.} Copy avoidance also depends on how the stack represents shared
data. For example, Linux copies \msgzc data delivered to local packet sockets, so a packet
tap can reintroduce copies~\cite{msgzerocopy}. Such behavior must be checked separately from
pinning and buffer management (\S\ref{sec:eval:tx}); it is not a universal property of
application-owned memory.
